\subsection*{発展的な読み物}
\addcontentsline{toc}{subsection}{9 発展的な読み物}
SWEBOK 第8章で紹介される発展的な読み物を、日本語で読む目的に合わせて整理する。

\par\smallskip\noindent\refstepcounter{table}\label{tab:ch08-09}表 \thetable: 発展的な読み物における文献・読むと得られること\par\smallskip
表 \thetable は、発展的な読み物における文献・読むと得られることを比較・確認しやすい形で整理したものである。\par\smallskip
\begin{longtable}{p{0.30\linewidth}p{0.64\linewidth}}\toprule
文献 & 読むと得られること \\ \midrule
\endfirsthead\toprule
文献 & 読むと得られること \\ \midrule
\endhead
Berczuk and Appleton, Software Configuration Management Patterns & 構成管理の実践を「パターン」として学べる。特定ツールに依存せず、チーム作業、統合、分岐、変更管理を考える助けになる。 \\
CMMI for Development & 成熟度レベル 2 の構成管理活動を含む。組織のプロセス改善と構成管理を結びつけて理解できる。 \\
Aiello and Sachs, Configuration Management Best Practices & 現実の組織で使える構成管理の実践を扱う。変更統制の種類や運用上の注意点を深く学べる。 \\
\bottomrule\end{longtable}
